% Scénarios de test consolidés — chapitre 3 (authentification, utilisateurs, sites)
\small
\PFETableSetup
\setlength{\LTpre}{0pt}
\setlength{\LTpost}{0pt}
\setlength{\LTleft}{-0.8cm}
\setlength{\LTright}{-0.8cm}
\renewcommand{\arraystretch}{1.25}

\begin{longtable}{|>{\centering\arraybackslash}m{3.1cm}|
>{\raggedright\arraybackslash}p{2.9cm}|
>{\raggedright\arraybackslash}p{4.4cm}|
>{\raggedright\arraybackslash}p{5.0cm}|
>{\centering\arraybackslash}p{1.3cm}|}
\caption{Synthèse des cas de test — Authentification \& Architecture de base}\label{tab:ch3-scenarios-tests}\\
\hline
\tableHeaderStyle
\tableHeaderCell{Nom du cas d'utilisation} &
\tableHeaderCell{Cas de test} &
\tableHeaderCell{Description du cas} &
\tableHeaderCell{Résultat attendu} &
\tableHeaderCell{Résultat réel} \\
\hline
\endfirsthead
\hline
\multicolumn{5}{|c|}{\tablename\ \thetable{} --- \textit{(suite)}} \\
\hline
\tableHeaderStyle
\tableHeaderCell{Nom du cas d'utilisation} &
\tableHeaderCell{Cas de test} &
\tableHeaderCell{Description du cas} &
\tableHeaderCell{Résultat attendu} &
\tableHeaderCell{Résultat réel} \\
\hline
\endhead
\hline
\multicolumn{5}{|r|}{\textit{Suite page suivante\ldots}} \\
\hline
\endfoot
\hline
\endlastfoot

\multirow{2}{=}{\centering S'authentifier} &
Identifiants manquants &
Tentative de connexion sans renseigner l'email ou le mot de passe. &
Message indiquant que les champs d'authentification sont obligatoires. &
Vrai \\
\cline{2-5}

 &
Identifiants invalides &
Tentative de connexion avec des informations incorrectes. &
Message d'échec d'authentification. &
Vrai \\
\hline

Accéder à Dashboard &
Utilisateur non authentifié &
L'utilisateur tente d'accéder au tableau de bord sans session valide. &
Redirection vers la page de connexion ou message d'accès refusé. &
Vrai \\
\hline

\multirow{5}{=}{\centering Gérer \\les utilisateurs} &
Afficher la liste des utilisateurs &
Ouverture du module de gestion des utilisateurs par un profil autorisé. &
La liste des utilisateurs s'affiche correctement. &
Vrai \\
\cline{2-5}

 &
Créer un utilisateur &
Création d'un nouveau compte avec des informations valides. &
Le compte est créé et apparaît dans la liste. &
Vrai \\
\cline{2-5}

 &
Modifier un utilisateur &
Mise à jour des informations d'un utilisateur existant. &
Les modifications sont enregistrées et visibles. &
Vrai \\
\cline{2-5}

 &
Désactiver un utilisateur &
Changement du statut d'un utilisateur actif vers inactif. &
Le compte est désactivé et l'état est mis à jour dans la liste. &
Vrai \\
\cline{2-5}

 &
Associer un utilisateur à un site avec un grade &
Attribution d'un utilisateur à un site avec un niveau de grade. &
L'affectation et le grade sont enregistrés correctement. &
Vrai \\
\hline

\multirow{4}{=}{\centering Gérer les sites} &
Créer un site &
Création d'un site avec les informations requises. &
Le site est créé et visible dans la liste des sites. &
Vrai \\
\cline{2-5}

 &
Modifier un site &
Modification des informations d'un site existant. &
Les nouvelles informations du site sont enregistrées. &
Vrai \\
\cline{2-5}

 &
Désactiver un site &
Passage d'un site actif en état désactivé. &
Le site est désactivé et n'est plus proposé aux actions actives. &
Vrai \\
\cline{2-5}

 &
Consulter un site &
Consultation des détails d'un site depuis la liste. &
Les informations détaillées du site sont affichées. &
Vrai \\
\hline

\multirow{4}{=}{\centering Gérer \\les documents} &
Télécharger un fichier &
Ajout d'un document dans la plateforme. &
Le fichier est importé et disponible dans la liste des documents. &
Vrai \\
\cline{2-5}

 &
Modifier un fichier &
Mise à jour des informations ou métadonnées d'un document existant. &
Les modifications sont enregistrées et visibles. &
Vrai \\
\cline{2-5}

 &
Supprimer un fichier &
Suppression d'un document depuis la gestion documentaire. &
Le fichier est supprimé et n'apparaît plus dans la liste. &
Vrai \\
\cline{2-5}

 &
Gérer les documents &
Accès global aux actions de consultation et de gestion documentaire. &
Les actions autorisées sont accessibles selon le profil utilisateur. &
Vrai \\
\hline

\end{longtable}
\normalsize
